{\small\opmaakproza
\begin{enumerate}[leftmargin=1.5em,itemsep=3pt,label=\textcolor{kleur1}{\sffamily\bfseries\arabic*}]
\item \textbf{Lees de inleiding en de afbeeldingen goed.} Biologie-examens geven veel nieuwe informatie in tekst,
  figuren en tabellen. Het antwoord staat vaak half in die informatie; jij koppelt het aan wat je weet.
\item \textbf{Geef zoveel antwoorden als gevraagd.} Vraagt de vraag om twee factoren, dan tellen alleen de eerste twee.
  Extra antwoorden kunnen je punten kosten als er een fout tussen zit.
\item \textbf{Gebruik de juiste biologische begrippen.} Schrijf \emph{osmose} in plaats van \emph{water trekt erin},
  en noem het orgaan, het weefsel of de stof bij zijn naam.
\item \textbf{Leg verbanden in stappen uit.} Een goede uitleg is een keten: dit verandert, daardoor gebeurt dat,
  en dus is het gevolg zo. Sla geen schakel over.
\item \textbf{Zoek op in je informatieboek.} Bouw van cellen, de genetische code, kringlopen en schema's van
  organen staan in Binas of ScienceData. Gebruik het, ook als je denkt het te weten.
\item \textbf{Lees grafieken en tabellen precies.} Kijk naar de assen, de eenheden en de legenda. Noem bij een
  conclusie de waarden waarop je hem baseert.
\item \textbf{Onderzoeksvragen}: let op variabelen, een controlegroep, herhaling en wat een proef wel en niet bewijst.
  Een conclusie gaat nooit verder dan de resultaten.
\item \textbf{Meerkeuzevragen}: lees alle alternatieven voordat je kiest, en streep af wat zeker niet klopt.
\item \textbf{Kijk naar de punten.} Een vraag van 2 punten vraagt meestal twee elementen. Houd ongeveer
  1,5 minuut per punt aan en ga door als je vastloopt.
\item \textbf{Blijf bij de vraag.} Beantwoord precies wat er staat. Een algemeen verhaal dat klopt maar niet aansluit
  bij de vraag, levert geen punten op.
\end{enumerate}\par}
